Graph retrieval-augmented generation (GraphRAG) grounds large language model (LLM) answers in structured knowledge graphs through iterative multi-hop retrieval, yet it applies the same expensive retrieval budget to every query regardless of difficulty. Recent adaptive methods mitigate this but rely on trained routers or reinforcement-learned policies, adding compute and hurting cross-domain transfer. We propose \method{}, a \emph{training-free, query-adaptive retrieval controller} that decides, per query and per hop, whether to use the graph at all, how wide the beam should be, and when to stop expanding---using only distributions already computed by the pipeline (entity-linking confidences and LLM relevance scores). \method{} shares a single retrieval engine with the fixed-budget baseline, so the controller is the only variable in our ablations by construction. On MetaQA with a 7B generator, \method{} strictly Pareto-dominates the fixed ToG-style baseline at 2-hop reasoning: it simultaneously improves F1 by $10.3\%$ (0.343$\rightarrow$0.378), Hit@1 by $7.0\%$, and reduces retrieval cost by $23.2\%$ (1111$\rightarrow$853 input tokens/query), a gap that is statistically significant ($p{=}0.049$, paired bootstrap). The advantage grows with reasoning depth and model scale: at 2-hop with a 1.5B model, F1 improves $23.5\%$ ($p{=}0.004$) at $19\%$ lower cost. Ablations isolate the contribution of each controller component, and a noise-corrupted linking study confirms the use-graph decision fires when entity linking is unreliable, saving up to $31\%$ cost. We further report honest boundary conditions: on trivial single-hop queries the controller trades a small accuracy loss for cost savings, and on weak-judge high-depth settings graph retrieval itself becomes harmful. \method{} requires no training, no learned parameters, and runs on a single 32\,GB GPU.
